%% Terms, identities, equational deduction, and free algebras.

\subsection{Terms, identities, and equational deduction}
\label{sec:p10-equational-logic}

\begin{hbdefinition}{def:p10-term-algebra}{Term algebra and evaluation}
Fix a finitary signature \(\tau\). For a set \(X\) disjoint from the
operation symbols, \(T_\tau(X)\) is the set of finite terms generated
by \(X\) and the symbols of \(\tau\) with their declared arities.
An \(n\)-ary basic operation on this set sends
\((t_1,\ldots,t_n)\) to the term \(f(t_1,\ldots,t_n)\).
In particular, a nullary symbol is itself a term.
Under the nonempty-algebra convention, this term algebra exists when
\(X\ne\varnothing\) or \(\tau\) has a nullary symbol.
If neither condition holds, there are no terms and the empty set is
not an algebra in the present convention.

For a \(\tau\)-algebra \(A\) and a map \(v:X\to A\), evaluation is
defined recursively by
\[
 x^A(v)=v(x),\qquad
 f(t_1,\ldots,t_n)^A(v)
       =f^A(t_1^A(v),\ldots,t_n^A(v)).
\]
A substitution \(\sigma:X\to T_\tau(Y)\) replaces each variable
simultaneously and extends recursively to a map
\(\widehat\sigma:T_\tau(X)\to T_\tau(Y)\).
Write \(t\sigma=\widehat\sigma(t)\).
No variable is bound in a term.
These constructions use the term-algebra presentation of
\cite[II\S10, Definitions~10.1--10.4]{BurrisSankappanavar2012}.
\index{term algebra}\index{substitution!equational}
\end{hbdefinition}

\begin{hbtheorem}{thm:p10-term-universal}{proved=annex:p10.eq.1}{Universal property and substitution law for terms}
If \(T_\tau(X)\) is nonempty, every \(v:X\to A\) extends to a unique
homomorphism \(T_\tau(X)\to A\), namely term evaluation.
For \(\sigma:X\to T_\tau(Y)\) and \(v:Y\to A\),
\[
 (t\sigma)^A(v)=t^A(w),
 \qquad w(x)=(\sigma(x))^A(v).
\]
Every algebra homomorphism \(h:A\to B\) also satisfies
\[
 h(t^A(v))=t^B(h\circ v).
\]
\end{hbtheorem}

\begin{hbdefinition}{def:p10-equational-logic}{Equational consequence}
Fix a variable set \(V\) for which \(T_\tau(V)\) is nonempty.
An equation \(s\approx t\) is an \emph{identity of \(A\)} when
\[
 A\models s\approx t
 \quad\Longleftrightarrow\quad
 \forall v:V\to A\quad s^A(v)=t^A(v).
\]
Write \(A\models\Gamma\) when every equation in \(\Gamma\) is an
identity of \(A\). Equational consequence is
\[
 \Gamma\models_{\mathrm{alg}}s\approx t
 \quad\Longleftrightarrow\quad
 \forall A\,
 (A\models\Gamma\ \Longrightarrow\ A\models s\approx t),
\]
where \(A\) ranges over nonempty \(\tau\)-algebras.
Thus each premise is required under every valuation.

The judgment \(\Gamma\vdash_{\mathrm{eq}}s\approx t\) means that a
finite derivation uses equations from \(\Gamma\), reflexivity,
symmetry, transitivity, compatibility with the basic operations,
and simultaneous term substitution. The complete rule forms are on
Sheet~\ref{sheet:equational-logic}.
Compatibility is equivalent to replacing one occurrence of a term
by an equal term inside a larger term; the proof below establishes
this correspondence with the source's replacement rule.
The semantic and deductive presentations are those of
\cite[II\S11, Definition~11.1; II\S14, Definitions~14.11,
14.16, 14.18]{BurrisSankappanavar2012}.
\index{equational logic}\index{identity!algebraic}
\end{hbdefinition}

\showcalculus{equational-logic}

\begin{hbexample}{def:p10-identities-valuations}{An identity and a local equation}
In the empty signature, the identity \(x\approx y\), for distinct
variables \(x,y\), requires an algebra to have exactly one element.
Consequently it implies every identity.
In contrast, the single valuation \(v(x)=v(y)=v(u)=0\), \(v(z)=1\)
in the set \(\{0,1\}\) satisfies \(x\approx y\) and fails
\(u\approx z\), where the displayed variables are distinct.
This local implication is a different consequence relation:
unrestricted substitution of a fixed local premise would not preserve it.
All uses of \(\models_{\mathrm{alg}}\) below have the universal
identity semantics just defined.
\end{hbexample}

\begin{hbdefinition}{def:p10-deductive-congruence}{The congruence of provable identities}
On \(T_\tau(V)\), put
\[
 s\mathrel{\theta_\Gamma}t
 \quad\Longleftrightarrow\quad
 \Gamma\vdash_{\mathrm{eq}}s\approx t.
\]
A congruence \(\theta\) on a term algebra is \emph{fully invariant}
when every endomorphism preserves it:
\[
 s\mathrel\theta t\ \Longrightarrow\
 \widehat\sigma(s)\mathrel\theta\widehat\sigma(t)
 \quad\text{for every }\sigma:V\to T_\tau(V).
\]
The deduction rules make \(\theta_\Gamma\) the least fully invariant
congruence containing the pairs of \(\Gamma\).
The exact leastness argument and the quotient-model construction are
given in the following proof
\cite[II\S14, Theorems~14.12 and~14.17]{BurrisSankappanavar2012}.
\index{congruence!fully invariant}
\end{hbdefinition}

\begin{hbtheorem}{thm:p10-equational-completeness}{proved=annex:p10.eq.2}{Completeness of equational logic}
For every set \(\Gamma\) of equations over \(V\) and every equation
\(s\approx t\) over \(V\),
\[
 \Gamma\models_{\mathrm{alg}}s\approx t
 \quad\Longleftrightarrow\quad
 \Gamma\vdash_{\mathrm{eq}}s\approx t.
\]
The nonempty quotient algebra \(T_\tau(V)/\theta_\Gamma\) satisfies
every identity in \(\Gamma\).
If the displayed consequence fails, the valuation \(x\mapsto[x]\)
in this quotient falsifies \(s\approx t\).
Moreover, consequence is finitary: for each consequence \(s\approx t\),
some finite \(\Gamma_0\subseteq\Gamma\) already implies \(s\approx t\).
\end{hbtheorem}

This is Birkhoff's completeness theorem, with a direct quotient proof
below; the source's deductive-closure proof is
\cite[II\S14, Theorem~14.19, p.~95]{BurrisSankappanavar2012}.
The locators in this section refer to the 2012 update of the
Millennium Edition, whose main-text pagination was aligned with the
1981 edition.

\subsection{Free algebras in an equational class}
\label{sec:p10-free-algebras}

\begin{hbdefinition}{def:p10-free-algebra}{Free algebra on a set of generators}
For a variety \(\mathcal V\), a free algebra on \(X\) is a pair
\((F_{\mathcal V}(X),\eta_X)\), where
\(F_{\mathcal V}(X)\in\mathcal V\) and
\(\eta_X:X\to F_{\mathcal V}(X)\), such that every map \(v:X\to A\)
with \(A\in\mathcal V\) factors uniquely as
\[
 v=\bar v\circ\eta_X,\qquad
 \bar v:F_{\mathcal V}(X)\longrightarrow A
 \quad\text{an algebra homomorphism}.
\]
The universal property is stated for the map \(\eta_X\);
injectivity is not part of this definition.
For the variety consisting of one-element algebras, all generators
have the same image.
\index{free algebra}\index{universal property}
\end{hbdefinition}

\begin{hbtheorem}{thm:p10-free-algebra}{proved=annex:p10.eq.3}{Construction and uniqueness of free algebras}
Let \(\mathcal V\) be the class of all nonempty models of a set
\(\Gamma\) of identities. If \(T_\tau(X)\ne\varnothing\), define
\[
 s\mathrel{\theta_{\mathcal V,X}}t
 \quad\Longleftrightarrow\quad
 \forall A\in\mathcal V\ \forall v:X\to A\quad
                    s^A(v)=t^A(v).
\]
Then
\[
 F_{\mathcal V}(X)=T_\tau(X)/\theta_{\mathcal V,X},
 \qquad \eta_X(x)=[x],
\]
is free in \(\mathcal V\) on \(X\).
Any two free algebras on \(X\) are uniquely isomorphic by an
isomorphism commuting with the generator maps
\cite[II\S10, Theorems~10.7, 10.10, 10.12;
II\S11, Theorem~11.4]{BurrisSankappanavar2012}.
\end{hbtheorem}

The hypothesis on \(T_\tau(X)\) matters when \(X=\varnothing\).
With a nullary operation, ground terms provide the construction.
Without nullary operations, the empty term set does not provide a
nonempty algebra. No uniform existence theorem for an empty set of
generators is then available under the nonempty-domain convention.
For example, algebras of the empty signature are nonempty sets.
An initial nonempty set \(F\) would have a unique map to \(\{0,1\}\),
but its two constant maps are distinct.
This example does not exclude special varieties that have an initial
algebra for other reasons.

For the signature consisting of a binary multiplication and a unit
constant, the monoid identities give the free monoid \(X^*\).
The empty word is the unit, concatenation is multiplication, and
\(\eta_X(x)\) is the one-letter word.
For \(v:X\to M\), the required extension sends
\(x_1\cdots x_n\) to \(v(x_1)\cdots v(x_n)\) and sends the empty
word to \(1_M\).
Associativity and the unit laws make it a homomorphism; every
homomorphism extending \(v\) must take these values.
In particular, the free monoid on the empty set is the one-element
monoid.
